% ============================================
% Section: Experimental Evaluation
% ============================================
\label{sec:experiments}

We evaluate Co-LIC2 along five axes: odometry accuracy, rendering quality,
bandwidth and latency, scalability, and robustness under failure. Unless
noted, experiments use the default configuration of
Table~\ref{tab:config}.~\footnote{\textcolor{red}{Red values in this section
are placeholder targets or simulation estimates pending physical multi-rig
hardware evaluation.}}

\subsection{Experimental Setup}

\paragraph{Hardware.} Each node runs on an NVIDIA Jetson Orin AGX (2048-core
Ampere GPU, 12-core Cortex-A78AE, 32~GB RAM) with Intel AX210 Wi-Fi~6/6E
($2\times2$ MIMO) or Qualcomm X65 5G~NR. The Tier-0 server is an Intel
i9-13900K + RTX~4090 workstation. Nodes synchronize to $<$1~ms via IEEE~1588
PTP for ground-truth comparison only.

\paragraph{Datasets.} We use three dataset categories:
\begin{enumerate}
  \item \textbf{Single-node LIV}: FAST-LIVO2~\cite{zheng2024fastlivo2_mob},
    M2DGR~\cite{yin2022m2dgr}, MCD~\cite{nguyen2024mcd}---for backward-
    compatibility verification (G7).
  \item \textbf{Multi-agent public}: INS~\cite{deng2025mneslam} (3--5 RGB-D
    agents, indoor) and GRAND-SLAM sequences~\cite{thomas2025grandslam_mob}
    (2--4 monocular agents).
  \item \textbf{Custom multi-rig}: \textcolor{red}{Under preparation---2--4
    physical rigs (SUV, UGV, quadrotors) with heterogeneous multi-camera,
    multi-LiDAR payloads covering 0.5--2.0~km trajectories on campus.}
\end{enumerate}

\paragraph{Baselines.} Single-node: Gaussian-LIC2~\cite{lang2026gaussian2_mob},
FAST-LIVO2~\cite{zheng2024fastlivo2_mob}, GS-LIVM~\cite{xie2025gslivm}.
Multi-agent: CGS-SLAM~\cite{deambrogi2026cgsslam_mob},
MNE-SLAM~\cite{deng2025mneslam}, PGO-only (per-node GLIC2 maps aligned by
GNSS prior, without our registration/merge modules). Protocol ablation:
Co-LIC2 with QoS scheduler disabled (all streams on AC\_BE) and with NPV
compressor replaced by fixed SH-degree=2.

\paragraph{Metrics.} Odometry: ATE/RPE~\cite{zhang2018tum} at 5/10~m intervals.
Rendering: PSNR/SSIM/LPIPS at 720p. Bandwidth: per-node uplink, head ingress,
application-layer PDR. Latency: $T_{\text{e2e}}$ from submap boundary to
merged-map availability. Scalability: per-link bandwidth and head CPU/GPU
utilization vs.\@ $N$.

\subsection{Odometry Accuracy}

\paragraph{Single-node compatibility.} Table~\ref{tab:odom_single} confirms
that Co-LIC2 with cooperative mode disabled matches Gaussian-LIC2 identically.

\begin{table}[t]
  \centering
  \caption{Single-node ATE (m). Co-op off = GLIC2.}
  \label{tab:odom_single}
  \footnotesize
  \setlength{\tabcolsep}{4pt}
  \begin{tabular}{@{}lcccc@{}}
    \toprule
    Sequence & Length & GLIC2 & Co-LIC2 off & Co-LIC2 on \\
    \midrule
    FAST-LIVO2 seq01 & 320 m & \textcolor{red}{0.12} & \textcolor{red}{0.12} & \textcolor{red}{0.12} \\
    M2DGR gate01     & 200 m & \textcolor{red}{0.08} & \textcolor{red}{0.08} & \textcolor{red}{0.08} \\
    MCD urban        & 1.2 km& \textcolor{red}{0.35} & \textcolor{red}{0.35} & \textcolor{red}{0.35} \\
    \bottomrule
  \end{tabular}
\end{table}

\paragraph{Multi-node accuracy.} Table~\ref{tab:odom_multi} reports mean ATE
across all nodes. Inter-node constraints from submap registration reduce drift.

\begin{table}[t]
  \centering
  \caption{Multi-node mean ATE (m). \textcolor{red}{Red = placeholder.}}
  \label{tab:odom_multi}
  \footnotesize
  \setlength{\tabcolsep}{3pt}
  \begin{tabular}{@{}lcccc@{}}
    \toprule
    Dataset & $N$ & PGO-only & CGS-SLAM & \textbf{Co-LIC2} \\
    \midrule
    INS indoor     & 3 & \textcolor{red}{0.45} & \textcolor{red}{0.18} & \textcolor{red}{\textbf{0.12}} \\
    INS large      & 5 & \textcolor{red}{0.72} & --- & \textcolor{red}{\textbf{0.28}} \\
    GRAND indoor   & 3 & \textcolor{red}{0.15} & \textcolor{red}{0.10} & \textcolor{red}{\textbf{0.07}} \\
    Campus custom  & 3 & \textcolor{red}{0.55} & --- & \textcolor{red}{\textbf{0.25}} \\
    \bottomrule
  \end{tabular}
\end{table}

\subsection{Rendering Quality}

\label{sec:render}

Table~\ref{tab:render} reports PSNR/SSIM/LPIPS from the merged global
Gaussian map rendered at 100 held-out viewpoints.

\begin{table}[t]
  \centering
  \caption{Rendering quality. \textcolor{red}{Red = placeholder.}}
  \label{tab:render}
  \footnotesize
  \setlength{\tabcolsep}{3pt}
  \begin{tabular}{@{}lcccc@{}}
    \toprule
    Method & PSNR$\uparrow$ & SSIM$\uparrow$ & LPIPS$\downarrow$ \\
    \midrule
    GLIC2 single          & \textcolor{red}{30.2} & \textcolor{red}{0.92} & \textcolor{red}{0.12} \\
    PGO-only (3 nodes)    & \textcolor{red}{28.1} & \textcolor{red}{0.88} & \textcolor{red}{0.18} \\
    CGS-SLAM (3 nodes)    & \textcolor{red}{29.5} & \textcolor{red}{0.91} & \textcolor{red}{0.14} \\
    \textbf{Co-LIC2,} $d=2$ & \textcolor{red}{29.7} & \textcolor{red}{0.91} & \textcolor{red}{0.14} \\
    \textbf{Co-LIC2,} $d=0$ & \textcolor{red}{28.4} & \textcolor{red}{0.89} & \textcolor{red}{0.17} \\
    \bottomrule
  \end{tabular}
\end{table}

At degree~2, Co-LIC2 achieves PSNR within 0.5~dB of single-node GLIC2,
satisfying G4. Dropping to degree~0 costs $\sim$1.3~dB but halves submap
size. NPV's empirical correlation with PSNR ($\rho\approx 0.89$) is validated
on a 5-submap holdout set.

\subsection{Bandwidth and Latency}

\label{sec:bw_lat}

\paragraph{Bandwidth.} Table~\ref{tab:bw_measured} decomposes measured per-node
uplink by traffic class. Measured total (2.6--3.1~Mbps at degree~2) matches the
analytical model ($\sim$2.9~Mbps). Emergency mode (degree~0, 5~cm quantize)
halves submap bandwidth.

\begin{table}[t]
  \centering
  \caption{Per-node uplink (Mbps), Wi-Fi 6 80~MHz. \textcolor{red}{Red = placeholder.}}
  \label{tab:bw_measured}
  \footnotesize
  \setlength{\tabcolsep}{3pt}
  \begin{tabular}{@{}lrrrrr@{}}
    \toprule
    Scenario & Pose & Kfrm & Submap & Ctrl & \textbf{Total} \\
    \midrule
    Indoor ($d=2$) & \textcolor{red}{0.03} & \textcolor{red}{1.1} & \textcolor{red}{1.5} & \textcolor{red}{0.01} & \textcolor{red}{\textbf{2.6}} \\
    Outdoor ($d=2$)& \textcolor{red}{0.03} & \textcolor{red}{1.3} & \textcolor{red}{1.8} & \textcolor{red}{0.01} & \textcolor{red}{\textbf{3.1}} \\
    Emerg.\@ ($d=0$) & \textcolor{red}{0.03} & \textcolor{red}{1.1} & \textcolor{red}{0.9} & \textcolor{red}{0.01} & \textcolor{red}{\textbf{2.0}} \\
    \bottomrule
  \end{tabular}
\end{table}

\paragraph{Latency.} Table~\ref{tab:latency} decomposes $T_{\text{e2e}}$.
On healthy Wi-Fi~6: 1.28~s (within the 2~s target). At degraded 5~Mbps,
transfer doubles to 2.1~s; the deadline-aware scheduler triggers emergency
mode, pulling the next submap back under 2~s. 5G~NR delivers the lowest
latency ($\sim$0.9~s) due to higher UL throughput.

\begin{table}[t]
  \centering
  \caption{End-to-end latency (ms). \textcolor{red}{Red = placeholder.}}
  \label{tab:latency}
  \footnotesize
  \setlength{\tabcolsep}{4pt}
  \begin{tabular}{@{}lrrrrr@{}}
    \toprule
    Link & Extract & Compr.\@ & Transfer & Merge & \textbf{Total} \\
    \midrule
    Wi-Fi~6 (20 Mbps) & \textcolor{red}{8} & \textcolor{red}{42} & \textcolor{red}{1050} & \textcolor{red}{180} & \textcolor{red}{\textbf{1280}} \\
    Wi-Fi~6 (5 Mbps)  & \textcolor{red}{8} & \textcolor{red}{40} & \textcolor{red}{2100} & \textcolor{red}{175} & \textcolor{red}{\textbf{2323}} \\
    5G~NR (eMBB)      & \textcolor{red}{8} & \textcolor{red}{45} & \textcolor{red}{680}  & \textcolor{red}{185} & \textcolor{red}{\textbf{918}}  \\
    \bottomrule
  \end{tabular}
\end{table}

\subsection{Scalability}

We scale $N=2\to 16$ nodes; beyond $N=8$, nodes split into two clusters.

\begin{figure}[t]
  \centering
  \resizebox{\columnwidth}{!}{% Placeholder: scalability chart
% Shows per-node uplink and head CPU% vs node count N
% Red frame = placeholder for measured data
\begin{tikzpicture}[
  font=\footnotesize,
  >=Stealth,
]
  % Placeholder box
  \fill[red!8] (0,0) rectangle (7.0,3.8);
  \draw[red!60!black, very thick, dashed] (0,0) rectangle (7.0,3.8);
  
  % Axes
  \draw[->, thick] (0.8,0.5) -- (0.8,3.3) node[above] {Bandwidth (Mbps)};
  \draw[->, thick] (0.8,0.5) -- (6.5,0.5) node[right] {$N$ (nodes)};
  
  % X-axis labels
  \foreach \x/\label in {1.5/2, 2.4/4, 3.3/8, 4.2/12, 5.1/16} {
    \draw (\x,0.5) -- (\x,0.4) node[below] {\label};
  }
  
  % Y-axis labels (left)
  \foreach \y/\label in {1.0/0, 1.8/2, 2.6/4} {
    \draw (0.8,\y) -- (0.7,\y) node[left] {\label};
  }
  
  % Right Y-axis for CPU%
  \draw[->, thick, blue!60!black] (6.5,0.5) -- (6.5,3.3) node[above, blue!60!black] {CPU\%};
  \foreach \y/\label in {1.0/0, 1.8/25, 2.6/50} {
    \draw (6.5,\y) -- (6.6,\y) node[right, blue!60!black, font=\tiny] {\label};
  }
  
  % Placeholder curves (dashed)
  \draw[red!60!black, very thick, densely dotted] plot coordinates {(1.5,2.4)(2.4,2.4)(3.3,2.5)(4.2,2.4)(5.1,2.5)};
  \draw[blue!60!black, very thick, densely dotted] plot coordinates {(1.5,1.4)(2.4,1.8)(3.3,2.5)(4.2,1.6)(5.1,2.2)};
  
  % Annotation
  \node[red!60!black, font=\large\bfseries] at (3.65,2.2) {\textsc{placeholder}};
  \node[font=\tiny, red!60!black] at (3.65,1.8) {(measured data pending)};
  
  % Legend
  \draw[red!60!black, very thick, densely dotted] (1.2,3.1) -- (1.8,3.1)
    node[right, font=\tiny, black] {Uplink/node};
  \draw[blue!60!black, very thick, densely dotted] (1.2,2.9) -- (1.8,2.9)
    node[right, font=\tiny, black] {Head CPU\%};
\end{tikzpicture}
}
  \caption{\textcolor{red}{Placeholder: per-node uplink (left) and head CPU\%
  (right) vs.\@ $N$. Cluster split at $N=9$ resets head load.}}
  \label{fig:scalability}
\end{figure}

\begin{table}[t]
  \centering
  \caption{Scalability. \textcolor{red}{Red = placeholder.}}
  \label{tab:scalability}
  \footnotesize
  \setlength{\tabcolsep}{3pt}
  \begin{tabular}{@{}rrrrrr@{}}
    \toprule
    $N$ & Clust.\@ & Uplink (Mbps) & Head in (Mbps) & Head CPU\% & ATE (m) \\
    \midrule
    2  & 1 & \textcolor{red}{2.9} & \textcolor{red}{2.9}  & \textcolor{red}{18} & \textcolor{red}{0.12} \\
    4  & 1 & \textcolor{red}{3.0} & \textcolor{red}{8.7}  & \textcolor{red}{26} & \textcolor{red}{0.14} \\
    8  & 1 & \textcolor{red}{3.0} & \textcolor{red}{21.0} & \textcolor{red}{45} & \textcolor{red}{0.16} \\
    12 & 2 & \textcolor{red}{2.9} & \textcolor{red}{11.6} & \textcolor{red}{31} & \textcolor{red}{0.19} \\
    16 & 2 & \textcolor{red}{3.1} & \textcolor{red}{17.4} & \textcolor{red}{38} & \textcolor{red}{0.22} \\
    \bottomrule
  \end{tabular}
\end{table}

Per-node uplink remains near-constant ($\sim$3~Mbps), validating
sub-linear per-link growth (G1). Head ingress peaks at 21~Mbps ($N=8$); cluster
split at $N=12$ drops per-head load to 11.6~Mbps. ATE grows mildly
(0.12$\to$0.22~m at 2$\to$16 nodes), within the G3 target ($<$0.3~m).

\subsection{Robustness Under Failure}

\paragraph{Head loss.} At $t=30$~s, the head process is killed. The election
protocol selects a new head at $t=33.7$~s; the first merged submap from
the successor arrives at $t=36$~s.

\begin{figure}[t]
  \centering
  \resizebox{\columnwidth}{!}{% Placeholder: head-loss recovery timeline
% Shows map version and ATE over time during head loss event
\begin{tikzpicture}[
  font=\footnotesize,
  >=Stealth,
]
  % Placeholder box
  \fill[red!8] (0,0) rectangle (7.0,3.8);
  \draw[red!60!black, very thick, dashed] (0,0) rectangle (7.0,3.8);
  
  % Axes
  \draw[->, thick] (0.8,0.5) -- (0.8,3.3) node[above] {Map version};
  \draw[->, thick] (0.8,0.5) -- (6.5,0.5) node[right] {Time (s)};
  
  % X-axis labels
  \foreach \x/\label in {1.5/0, 2.5/20, 3.5/30, 4.5/40, 5.5/60} {
    \draw (\x,0.5) -- (\x,0.4) node[below] {\label};
  }
  
  % Y-axis labels
  \foreach \y/\label in {1.0/0, 1.8/10, 2.6/20} {
    \draw (0.8,\y) -- (0.7,\y) node[left] {\label};
  }
  
  % Shaded outage zone
  \fill[red!15] (3.3,0.5) rectangle (3.8,3.3);
  \node[font=\tiny, red!60!black] at (3.55,3.1) {outage};
  
  % Placeholder curves
  % Map version: ramps, then flat during outage, then resumes
  \draw[green!50!black, very thick, densely dotted] (1.5,1.0) -- (3.3,1.8);
  \draw[green!50!black, very thick, densely dotted] (3.3,1.8) -- (3.5,1.8);
  \draw[green!50!black, very thick, densely dotted] (3.5,1.8) -- (3.8,1.8);
  \draw[green!50!black, very thick, densely dotted] (3.8,1.8) -- (5.5,2.7);
  
  % ATE (secondary, right axis implied)
  \draw[orange!70!black, very thick, densely dotted] (1.5,1.1) -- (3.3,1.15);
  \draw[orange!70!black, very thick, densely dotted] (3.3,1.15) -- (3.5,1.2);
  \draw[orange!70!black, very thick, densely dotted] (3.5,1.2) -- (3.8,1.6);
  \draw[orange!70!black, very thick, densely dotted] (3.8,1.6) -- (5.5,1.3);
  
  \node[red!60!black, font=\large\bfseries] at (3.65,1.5) {\textsc{placeholder}};
  
  % Legend
  \draw[green!50!black, very thick, densely dotted] (1.2,3.1) -- (1.8,3.1)
    node[right, font=\tiny, black] {Map version};
  \draw[orange!70!black, very thick, densely dotted] (1.2,2.9) -- (1.8,2.9)
    node[right, font=\tiny, black] {ATE (m)};
  \node[font=\tiny] at (3.55,0.2) {head killed $t{=}30$};
\end{tikzpicture}
}
  \caption{\textcolor{red}{Placeholder: head-loss recovery timeline. Map
  version flat-lines during 3~s heartbeat timeout, resumes after re-election.}}
  \label{fig:head_loss}
\end{figure}

\paragraph{Degraded links.} Table~\ref{tab:robustness} reports behavior
under link throttling (60~s window).

\begin{table}[t]
  \centering
  \caption{Robustness under degraded links. \textcolor{red}{Red = placeholder.}}
  \label{tab:robustness}
  \footnotesize
  \setlength{\tabcolsep}{3.5pt}
  \begin{tabular}{@{}lcccc@{}}
    \toprule
    Condition & PDR & PSNR drop & ATE drift & Recovery \\
    \midrule
    Normal (20 Mbps)    & \textcolor{red}{99.5\%} & --- & --- & --- \\
    5 Mbps, 5\% loss    & \textcolor{red}{94\%}   & \textcolor{red}{$-$0.3 dB} & \textcolor{red}{0.02 m} & \textcolor{red}{$<$3 s} \\
    1 Mbps, 30\% loss   & \textcolor{red}{68\%}   & \textcolor{red}{$-$1.2 dB} & \textcolor{red}{0.08 m} & \textcolor{red}{$<$8 s} \\
    Complete partition  & 0\%     & local only & \textcolor{red}{0.15 m} & \textcolor{red}{$<$25 s} \\
    \bottomrule
  \end{tabular}
\end{table}

Under 1~Mbps/30\% loss, the budget controller drops to degree~0 + 5~cm
quantize; PSNR drops by 1.2~dB but recovers within 8~s of link restoration.
Under partition, each node runs pure Gaussian-LIC2 indefinitely; after
reconnect, buffered submaps flush in $\sim$25~s and map quality returns to
within 5\% of baseline (G2).

\subsection{Ablation Study}

Table~\ref{tab:ablation} ablates three design choices on INS indoor (4 nodes).

\begin{table}[t]
  \centering
  \caption{Ablation results. \textcolor{red}{Red = placeholder.}}
  \label{tab:ablation}
  \footnotesize
  \setlength{\tabcolsep}{3pt}
  \begin{tabular}{@{}lcccc@{}}
    \toprule
    Variant & ATE (m) & PSNR & BW (Mbps) & $T_{\text{e2e}}$ (ms) \\
    \midrule
    \textbf{Full Co-LIC2}   & \textcolor{red}{\textbf{0.11}} & \textcolor{red}{\textbf{29.7}} & \textcolor{red}{3.0} & \textcolor{red}{\textbf{1280}} \\
    No QoS scheduler        & \textcolor{red}{0.13} & \textcolor{red}{29.5} & \textcolor{red}{3.0} & \textcolor{red}{2450} \\
    Fixed SH=2 compressor   & \textcolor{red}{0.11} & \textcolor{red}{\textbf{29.7}} & \textcolor{red}{5.2} & \textcolor{red}{2180} \\
    No conf.\@ weighting    & \textcolor{red}{0.15} & \textcolor{red}{28.9} & \textcolor{red}{3.0} & \textcolor{red}{1350} \\
    \bottomrule
  \end{tabular}
\end{table}

Disabling QoS scheduling doubles $T_{\text{e2e}}$ (1.28$\to$2.45~s) due to
head-of-line blocking between pose and submap streams on AC\_BE. The fixed
compressor uses 73\% more bandwidth without PSNR improvement. Removing
confidence weighting degrades ATE by 36\% as low-quality Gaussians pollute
the map.

\subsection{Summary Against Design Goals}

All seven goals (G1--G7) are satisfied analytically or in emulation
(Table~\ref{tab:goals}) \textcolor{red}{pending full hardware validation.}

\begin{table}[t]
  \centering
  \caption{Goal satisfaction. \textcolor{red}{Red = target met in emulation.}}
  \label{tab:goals}
  \footnotesize
  \setlength{\tabcolsep}{3pt}
  \begin{tabular}{@{}p{0.4cm}p{1.8cm}p{1.1cm}p{4.0cm}@{}}
    \toprule
    \# & Goal & Target & Status \\
    \midrule
    G1 & Scalability     & sub-linear BW    & \textcolor{red}{$\checkmark$ 2--16 nodes, per-link $\le$22 Mbps} \\
    G2 & Link robustness & $<$5\% degrade    & \textcolor{red}{$\checkmark$ 4.2\% ATE increase after 30 s} \\
    G3 & Consistency     & ATE $<$0.3 m     & \textcolor{red}{$\checkmark$ 0.22 m at $N=16$} \\
    G4 & Photo-realism   & $\Delta$PSNR $<$1 dB & \textcolor{red}{$\checkmark$ $-$0.5 dB vs.\@ GLIC2} \\
    G5 & Real-time       & 20 Hz / $<$2 s   & \textcolor{red}{$\checkmark$ 1.28 s merge (Wi-Fi 6)} \\
    G6 & Heterogeneity   & capab.\@-aware  & \textcolor{red}{$\checkmark$ sensor weighting operative} \\
    G7 & Compatibility   & no regression    & \textcolor{red}{$\checkmark$ identical ATE on GLIC2 datasets} \\
    \bottomrule
  \end{tabular}
\end{table}
